\clearpage
\fancyhead[L]{\small NOVA ACTA UPSALIENSIS.}
\hypertarget{native-609-lower}{}
\begin{center}\small Göttingische gelehrte Anzeigen. Stück 70. Seiten 692..694. 2. Mai 1818.\\[3pt]\rule{.72\linewidth}{.3pt}\end{center}

\emph{Nova Acta regiae societatis scientiarum Upsaliensis.} Vol. VII. Upsala. 394 Seiten im Quartformat.

\begingroup\brokenpenalty=10000\relax
Diese Sammlung beginnt mit einer vom verstorbenen Prof. \textsc{Erich Prosperin} hinterlassenen Abhandlung über die Bahnen, die die Nebenplaneten um die Sonne beschreiben würden, wenn ihre Hauptplaneten plötzlich vernichtet würden oder aufhörten, auf sie einzuwirken.
Je nach den Umständen können diese Bahnen Ellipsen, Parabeln oder Hyperbeln werden, rechtläufig oder rückläufig.
Die Art des Kegelschnitts hängt teils von der Entfernung des Nebenplaneten von seinem Hauptplaneten ab, teils von der Stelle, an der sich der Nebenplanet bei der vorausgesetzten Vernichtung des Hauptplaneten gerade auf seiner Bahn befindet.
Wenn man mit dem Verfasser die Elliptizität der Bahn des Hauptplaneten um die Sonne und die der ursprünglichen Bahn des Nebenplaneten um seinen Hauptplaneten vernachlässigt und das Produkt aus der Entfernung des Hauptplaneten von der Sonne und seiner Masse (wobei die Masse der Sonne als Einheit genommen wird) $= d$ setzt, die Entfernung des Nebenplaneten von seinem Hauptplaneten hingegen $= D$, dann wird die neue Bahn des Nebenplaneten um die Sonne nach der Vernichtung des Hauptplaneten immer eine Ellipse sein, wenn $D$ größer ist als $(3+\BareRoot{8})d$.
Sie wird immer eine Hyperbel werden, wenn $D$ kleiner ist als $(3-\BareRoot{8})d$.
Zwischen beiden Grenzen kann sie eine Ellipse, eine Parabel oder eine Hyperbel werden.
Ist $D$ größer als $d$, so kann die neue Bahn nicht rückläufig werden.
Ist $D$ größer als $d\BareRoot{\tfrac{1}{2}}$, so kann sie zumindest keine rückläufige Hyperbel werden.
Ist schließlich $D$ kleiner als $d\BareRoot{\tfrac{1}{8}}$, so kann sie keine rechtläufige Ellipse bleiben.
Diese Sätze setzen voraus, dass die Masse des Hauptplaneten im Vergleich zu der der Sonne sehr klein ist und dass die ursprüngliche Bewegung des Nebenplaneten in der Ebene des Hauptplaneten stattfindet. Ohne diese Voraussetzungen bedürfen die Sätze einiger Änderungen.
Herr \textsc{Prosperin} hat jeden einzelnen Nebenplaneten unseres Sonnensystems gesondert betrachtet und sich die Mühe gemacht, dafür sehr umfangreiche Tabellen zu berechnen.
Die Bahn unseres Mondes würde, wo auch immer er die Erde plötzlich verlöre, eine rechtläufige Ellipse bleiben. Die Jupiter-, Saturn- und Uranustrabanten hingegen könnten je nach den Umständen auch parabolische oder hyperbolische Bahnen erhalten.
Offenbar lassen sich diese Sätze nun auch umkehren. \SourceLeaf{610}Ein Weltkörper, der auf einer parabolischen oder hyperbolischen Bahn in die Nähe des Jupiter, Saturn oder Uranus gelangte, würde eine kreisförmige Bahn um den Planeten antreten, wenn dessen Einwirkung erst nach bereits erfolgter Annäherung \emph{mit einem Male} anfinge.
Da es sich in der Natur aber anders verhält, bleibt die ganze Untersuchung eigentlich nichts als ein mathematisches Spielwerk.
Diesen Umstand hat der Verfasser zwar nicht ganz übersehen, aber nicht genügend beherzigt, denn es ist genau ebenso \emph{unmöglich}, dass ein Komet, der zufällig in die Nähe eines Hauptplaneten kommt, durch die bloßen Anziehungskräfte zu dessen Trabanten werden kann, wie dass ein bereits vorhandener Nebenplanet seinen Hauptplaneten verlässt, um sich auf einer parabolischen oder hyperbolischen Bahn zu entfernen.
\par\clearpage\endgroup
